\documentclass[../../../main.tex]{subfiles}

\begin{document}

\chapter{[N] Diophantine Equations}

Diophantine equation refers to equations with integer coefficient
whose interests are integer solutions. For instance, equations like
$ax + by = c$, $x^2 - ny^2 = 1$, and $x^n + y^n = z^n$ whose
interests lie only on integer solutions are Diophantine equations.
Diophantine equation has very wide range in difficulties starting
from finding Pythagorean triples to famous theorems like
Fermat's Last Theorem. In this note, we will discuss foundational
topics on Diophantine equations to common techniques like proof
by infinite descent and Vieta jumping.

\section{Primitive Pythagorean Triples}

Problems on Pythagorean triples are probably one of the first
problems that most of us solve in number theory, and yet
remains as one of the dominant theorem. Maybe this is slightly
biased towards Korean MO, but Pythagorean triples are very
common and worth noting. As always, let's start with foundational
properties.

\begin{theorem}{}{Euclid's Formula}
For a primitive Pythagorean triple $(a, b, c)$, there exists a pair
of positive integers $m > n$ such that the following equations hold.
\[
\{ a, b \} = \{ m^2 - n^2, 2mn \}, \quad
c = m^2 + n^2
\]
\end{theorem}

Before we prove this theorem, let's clarify few terms. A primitive
Pythagorean triple $(a, b, c)$ means triples $(a, b, c)$ such that
$(a, b) = (b, c) \linebreak = (c, a) = 1$. Also, we cannot
definitively say that $a = m^2 - n^2$ and $b = 2mn$ because it depends
on which term is odd and even. Indeed, we have opposite parities for
$a$ and $b$. I mean we can technically say $a = m^2 - n^2$ and
$b = 2mn$ without loss of generality, but let's keep our theorem
this way to highlight the parity. Now here is the famous proof!

\begin{proof}
Without loss of generality, let $2 \nmid a$. By construction,
the following congruences hold.
\begin{align*}
    a^2 \equiv 1 \pmod{4}&, \quad b^2 \equiv 0 \pmod{4} \\
    c^2 \equiv a^2 + b^2 &\equiv 1 \pmod{4}
\end{align*}
Therefore, $c$ is odd. Let $a = 2a' + 1$ and $c = 2c' + 1$.
Then, the following congruences hold.
\begin{align*}
    (2a' + 1)^2 + b^2 &\equiv (2c' + 1) \\
    4a'^2 + 4a' + 1 + b^2 &\equiv 4c'^2 + 4c' + 1 \pmod{8} \\
    b^2 \equiv 4c'^2 + 4c' - 4a'^2 - 4a'
    &\equiv 4 (c'^2 + c' - a'^2 - a') \equiv 0 \pmod{8}
\end{align*}
Thus $4 \mid b$. Continuing, the following equation holds by
definition.
\[
(c + a)(c - a) = b^2
\]
Let $d = (c + a, c - a)$. Note that $2 \mid d$ since $2 \nmid a, c$.
Moreover, $d \mid (c + a) + (c - a) = 2c$ and $d \mid (c + a) -
(c - a) = 2a$. Because $(a, c) = 1$, $d = 2$. In other words,
$c + a = 2m^2$ and $c - a = 2n^2$ for some coprime positive
integers $m, n$. Substituting, $b = 2mn$. Thus, there exists
a pair of positive integers $m > n$ such that $\{ a, b \} =
\{ m^2 - n^2, 2mn \}$ and $c = m^2 + n^2$.
\end{proof}

Honestly, this is pretty much it for the section. There isn't
lots of topics to cover for primitive Pythagorean triple, but
its applications are very worth noting. We will discuss more
on the application of the theorem as we discuss infinite
descent, but before we conclude this section, let's solve
a practice problem to see how it is used in problems.

\begin{important}
When using this theorem, it is very important to consider the
condition stated that $(a, b) = (b, c) = (c, a) = 1$. Moreover,
it is important to consider the parity of $a$ and $b$ before
assigning respective values $m^2 - n^2$ and $2mn$. Lastly, but
most importantly, remember that $(m, n) = 1$ and they are opposite
parity! The last condition is very easy to miss (at least for me)
and used very often when we apply this theorem recursively.
\end{important}

Below is an application that seems obvious, but is actually
way difficult than it looks. For some reasons, I got lost very
easily when learning applications of the theorem and any topics
in NT, and I keeping the conditions of given and new variables
in mind really helped me.

\begin{problem}{2020 Korean MO Problem 4}{2020 Korean MO Problem 4}
Find a pair of positive coprime integers $(a, b)$ other than
$(41, 12)$ such that $a^2 - 5b^2$ and $a^2 + 5b^2$ are both perfect
squares.
\end{problem}
This solution is largely based on
\href{https://youtu.be/4eetdspSRYw?si=U1hylcoUzUwPA66a}{a solution by WinnersKMO}.
\begin{solution}
First, let $(x, y)$ be another coprime pair such that
$x^2 - 5y^2 = m^2$ and $x^2 + 5y^2 = n^2$ for some integers
$m$ and $n$. Rearranging the equation, $2x^2 = m^2 + n^2$ and
$10y^2 = n^2 - m^2$ hold. Starting with the first equation,
the following equation is obtained.
\[
x^2
= \frac{m^2 + n^2}{2}
= \left( \frac{m + n}{2} \right)^2 + \left( \frac{m - n}{2} \right)^2
\]
From the equation, assume that $x$, $\frac{m + n}{2}$, and
$\frac{m - n}{2}$ are pairwise coprime. Note that we are allowed
to set this condition since we just need to find any $x$ and $y$.
By Theorem \ref{theo:Euclid's Formula}, there exists coprime
integers $u, v$ with opposite parity such that $x = u^2 + v^2$,
$\frac{m - n}{2} = u^2 - v^2$, and $\frac{m + n}{2} = 2uv$.
Again, we are allowed to set parities since we are finding any
$x$ and $y$.

Continuing, the second equation can be modified.
\begin{align*}
    10y^2
    &= 4 \cdot \frac{n + m}{2} \cdot \frac{n - m}{2} \\
    \frac{5}{2} y^2
    &= (u^2 - v^2) 2uv \\
    \therefore
    5 \left( \frac{y}{2} \right)^2
    &= (u + v)(u - v) uv
\end{align*}
By construction, $u, v$ are coprime to each other and $(u + v)(u - v)$
and $uv$ are also coprime. Let $u = 41^2$ and $v = 5 \cdot 12^2$.
Notice that the conditions that we have set while defining the
variables are met. Moreover as given by the problem, it is evident
that $(41^2 + 5 \cdot 12^2)(41^2 - 5 \cdot 12^2)$ are perfect
squares, returning an integer $y$. Hence, the following equations
are obtained.
\begin{align*}
    5 \left( \frac{y}{2} \right)^2
    &= (41^2 + 5 \cdot 12^2)(41^2 - 5 \cdot 12^2)
        \cdot 41^2 \cdot 5 \cdot 12^2 \\
    \left( \frac{y}{2} \right)^2
    &= (41^2 + 5 \cdot 12^2)(41^2 - 5 \cdot 12^2)
        \cdot 41^2 \cdot 12^2 \\
    \therefore y
    &= 2 \cdot 12 \cdot 41 \cdot
        \sqrt{(41^2 + 5 \cdot 12^2)(41^2 - 5 \cdot 12^2)}
    = 1494696
\end{align*}
Moreover, $x = u^2 + v^2 = \left(41^2\right)^2 +
\left(5 \cdot 12^2\right)^2 = 3344161$. Therefore,
$(3344161, 1494696)$ is another pair.
\end{solution}

With this in mind, let's discuss Pell's equation, another fundamental
topic from Diophantine equations.

\end{document}


